\chapter*{Аннотация}
\addcontentsline{toc}{chapter}{Аннотация}

Определили, что скорость откачки воздуха из вакуумной системы по зависимости давления от времени в процессе откачки при помощи диффузионного насоса составила $W = (3.17 \pm 0.13) \cdot 10^2~\text{см}^3/\text{с}$. Также нашли величину скорости откачки через капилляр после введения в систему течи $W = (27 \pm 10) \ \text{см}^3 / \text{с}$. Скорость откачки через узкий капилляр оказалась меньше, чем через трубку большего диаметра. Это показывает, что скорость откачки зависит от пропускной способности трубки. Сделан вывод, что после создания течи эффективная пропускная способность системы вместе с капилляром снизилась. Наличие течи не позволило достичь максимальной скорости откачки.